\section{Technical Details for the Main Bound}
\label{app:full-proof}

This appendix expands the normal forms, uniform constant checks, kernel
bookkeeping, and boundary cases used in \cref{sec:theory}.  It introduces no
new hypothesis or proof step.  As in the main proof, the frozen output is an
exact mathematical object; the neural model and experimental distribution play
no role.

\subsection{Final sequence and main theorem}

Recall the coprime integers

\[
  A=582942583375009,\qquad B=250000000000000,\qquad
  R=A/B>1
\]

with $\gcd(A,B)=1$ and $A>B\ge2$.  All asymptotic window and conductor
statements below are for $t\ge t_*$, where $t_*$ is a fixed index large enough
that $h_t\ge4$, $m(t)\ge2$, and $h_{t+1}>h_t+1$.  The finitely many smaller
indices are handled separately in the final summation.

Recall also the backbone $h_t=\lfloor\alpha R^t\rfloor$ from
\cref{eq:geometric-backbone}.  Define

\begin{equation}
  \kappa:=\log_A R
  =0.024901468995116044855\ldots.
  \label{eq:kappa}
\end{equation}

The practical sequence $\mathcal H_{\rm exec}$ is given by
\cref{eq:exec-sequence}.  To define its asymptotic completion, fix

\[
  N_{\rm guard}=10^{1000},\qquad s_0=2724,
\]

and recursively set

\begin{equation}
  s_{r+1}=s_r+\left\lceil\log(s_r+2)\right\rceil,
  \label{eq:companion-indices}
\end{equation}

where $\log$ is natural.  At every $s_r$, add the unit companion
$h_{s_r}+1$.  The final infinite set is

\begin{equation}
  \mathcal H^\sharp_{\rm tune}
  :=\mathcal H_{\rm exec}
    \cup\{h_{s_r}+1:r\ge0\}.
  \label{eq:completed-sequence}
\end{equation}

As usual, the Shellsort for length $N$ uses all distinct members of this set
below $N$, in decreasing order.

\begin{theorem}[Main asymptotic bound]
  \label{thm:main}
  The sequence $\mathcal H^\sharp_{\rm tune}$ has the following properties.
  \begin{enumerate}
    \item For every $N\le10^{1000}$, its gap list is exactly the list emitted
    by $\mathcal H_{\rm exec}$.  Hence all passes, comparisons, moves, and
    memory accesses are identical.
    \item The number of original gaps below $N$ is $\Theta(\log N)$, whereas
    the number of unit companions is
    \[
      \bigO\!\left(\frac{\log N}{\log\log N}\right)=o(\log N).
    \]
    \item Its worst-case operation count is
    \begin{equation}
      T_{\mathcal H^\sharp_{\rm tune}}(N)
      =\bigO\!\left(
        N^{1.024296451657\ldots}\polylog N
      \right).
      \label{eq:main-bound}
    \end{equation}
    \item Zang's rational-geometric lower bound~\cite[Theorem~3]{zang2026}
    applies to this completed sequence and gives
    \[
      T_{\mathcal H^\sharp_{\rm tune}}(N)
      =\Omega\!\left(N^{1.024296451657\ldots}\right).
    \]
    Hence the polynomial exponent in~\eqref{eq:main-bound} is tight, up to its
    polylogarithmic factor.
  \end{enumerate}
\end{theorem}

The completion in~\cref{eq:completed-sequence} is part of the theorem, not an
assumption about the learned backbone.  We do not claim that the uncompleted
geometric set alone satisfies~\eqref{eq:main-bound}.  The point of the guard
threshold is that the theorem and the finite experiment refer to exactly the
same visible sequence while the infinite tail is made arithmetically
certifiable.

\subsection{Pass cost as a multiplier-semigroup problem}

Fix a current gap $d$.  Before the $d$-pass, suppose the larger gaps
$a_1,\ldots,a_m$ have already been executed.  Define

\[
  M_d(a_1,\ldots,a_m)
  =\left\{k\in\Z_{\ge0}:
    kd=\sum_{i=1}^m c_i a_i,\quad c_i\in\Z_{\ge0}
  \right\}
\]

and let $U_d$ be the number of missing positive multipliers, with
$U_d=\infty$ if the complement is infinite.  Incerpi and Sedgewick's
generalized Frobenius pass lemma gives~\cite[Lemma~1,
pp.~212--213]{incerpi1985}

\begin{equation}
  \cost_d(N)=\bigO\bigl(N(U_d+1)\bigr).
  \label{eq:frobenius-pass}
\end{equation}

There is also an elementary bound.  The $d$ residue classes each have length
at most $\lceil N/d\rceil$, so insertion-sorting them costs

\begin{equation}
  \cost_d(N)=\bigO(N^2/d).
  \label{eq:trivial-pass}
\end{equation}

Combining them,

\begin{equation}
  \cost_d(N)=\bigO\!\left(
    \min\left\{N^2/d,\ N(U_d+1)\right\}
  \right).
  \label{eq:combined-pass}
\end{equation}

It remains to prove that all sufficiently large $k$ belong to $M_d$.  If the
least such threshold is $K$, then trivially $U_d\le K-1$.

\subsection{The logarithmic future window}

For a backbone gap $h_t$, take

\begin{equation}
  m(t):=\min\{m\in\Z_{\ge0}:A^m\ge h_t\}
       =\lceil\log_A h_t\rceil.
  \label{eq:window}
\end{equation}

Since $h_t=\alpha R^t+O(1)$,

\begin{equation}
  m(t)=\kappa t+O(1).
  \label{eq:window-index}
\end{equation}

For $d\asymp h_t$ and $m=m(t)$ we have

\begin{equation}
  A^m=\Theta(d),\qquad
  B^m=\Theta(d^{1-\kappa}),\qquad
  \frac d{B^m}=\Theta(d^\kappa),\qquad
  R^m=\Theta(d^\kappa).
  \label{eq:window-scales}
\end{equation}

Indeed, $A^{m-1}<h_t\le A^m$.  This remains true up to fixed multiplicative
constants for either possible tail gap $d\in\{h_t,h_t+1\}$.  Since
$R=A^\kappa$, we have $R^m=(A^m)^\kappa=\Theta(d^\kappa)$ and
$B^m=A^m/R^m=\Theta(d^{1-\kappa})$, proving all four estimates with constants
independent of $t$.

The increasing window is essential.  A fixed number of rational-geometric
future gaps retains a denominator-induced congruence obstruction whose modulus
grows with $d$; it cannot yield the sublinear missing-multiplier count needed
below.

\subsection{Exact denominator clearing}

Set $d=h_t$ for the moment and take the future backbone gaps

\[
  a_i=h_{t+i},\qquad1\le i\le m.
\]

Define

\begin{equation}
  E_i=B^i a_i-A^i d,\qquad
  q_i=A^iB^{m-i},\qquad
  v_i=B^{m-i}E_i.
  \label{eq:qv-definition}
\end{equation}

Then

\begin{equation}
  B^m a_i=dq_i+v_i.
  \label{eq:two-coordinate}
\end{equation}

Consequently $kd=\sum_i c_i a_i$ holds if and only if there is an integer
$D$ such that

\begin{equation}
  \sum_i c_iq_i=kB^m-D,\qquad
  \sum_i c_iv_i=Dd.
  \label{eq:coupled-coordinates}
\end{equation}

This is an exact two-coordinate reformulation.  The $q$-coordinate captures
the rational geometric scale, the $v$-coordinate captures flooring error, and
$D$ couples them.

For uniform error control, write

\[
  \theta_t=\alpha R^t-\lfloor\alpha R^t\rfloor\in[0,1).
\]

A direct calculation gives

\begin{equation}
  \sigma_i:=\frac{v_i}{q_i}
  =\frac{E_i}{A^i}
  =\theta_t-R^{-i}\theta_{t+i},
  \qquad |\sigma_i|<1.
  \label{eq:slopes}
\end{equation}

If the current gap is the auxiliary value $d=h_t+1$, every slope is translated
by $-1$; their width is unchanged and $|\sigma_i|<2$.  All constants below may
therefore depend on $A,B$ but not on $t,m,$ or $d$.

\subsection{Compound coordinates and their complete carry kernel}

Every $q_i$ contains a factor $A$.  Put

\[
  w_i=q_i/A=A^{i-1}B^{m-i}.
\]

Then

\[
  (w_1,\ldots,w_m)
  =(B^{m-1},AB^{m-2},\ldots,A^{m-1}),
  \qquad Aw_i=Bw_{i+1}.
\]

These are a constant-ratio compound sequence.  The following standard normal
form is a specialization of the compound numerical-semigroup theory of
Kiers, O'Neill, and Ponomarenko~\cite[Proposition~12, Theorem~15,
Corollary~16]{kiers2016compound}; we include the argument needed here.

\begin{lemma}[Compound threshold]
  \label{lem:compound}
  For the semigroup $\mathcal S_m=\langle w_1,\ldots,w_m\rangle$, every
  representable integer has a unique normal form
  \[
    X=\sum_{i=1}^m c_iw_i,\qquad0\le c_i<B\quad(2\le i\le m).
  \]
  Its Frobenius number is
  \begin{equation}
    F_m=(B-1)\sum_{i=2}^m w_i-w_1=\bigO(A^{m-1}).
    \label{eq:compound-frobenius}
  \end{equation}
  Hence every $X>F_m$ lies in $\mathcal S_m$.
\end{lemma}

\begin{proof}
  Starting at the largest index, divide $c_{i+1}$ by $B$.  Keep its remainder
  at coordinate $i+1$ and transfer $A$ times the quotient to coordinate $i$,
  using $Bw_{i+1}=Aw_i$.  This produces the claimed digit bounds.  Reading the
  equation successively modulo $B$ proves uniqueness.  Relative to the least
  generator $w_1=B^{m-1}$, the Ap\'ery set is exactly
  \[
    \left\{\sum_{i=2}^m c_iw_i:0\le c_i<B\right\}.
  \]
  Its largest member minus $w_1$ gives~\eqref{eq:compound-frobenius}.  Since
  the largest $w_i$ is $A^{m-1}$ and $A,B$ are fixed, the stated asymptotic
  bound follows.
\end{proof}

Now define adjacent carry vectors

\begin{equation}
  r_i=A\mathbf e_i-B\mathbf e_{i+1},\qquad1\le i<m.
  \label{eq:carry-vector}
\end{equation}

They preserve the main coordinate because $q\cdot r_i=0$.

\begin{proposition}[Complete integer carry kernel]
  \label{prop:kernel}
  The adjacent carries form a basis of the full integer kernel:
  \begin{equation}
    \ker_{\Z}(q)
    =\bigoplus_{i=1}^{m-1}\Z r_i.
    \label{eq:full-kernel}
  \end{equation}
\end{proposition}

\begin{proof}
  Since $q=Aw$, the two vectors have the same integer kernel.  Let
  $x\in\ker_{\Z}(w)$.  Reduce $w\cdot x=0$ modulo $A$.  Only
  $w_1=B^{m-1}$ is nonzero modulo $A$, and it is invertible because
  $\gcd(A,B)=1$; hence $A\mid x_1$.  Write $x_1=Az_1$ and subtract
  $z_1r_1$.  The first coordinate becomes zero.  Removing it and dividing the
  remaining weights by $A$ leaves the same problem of length $m-1$.  Induction
  yields $x=\sum_i z_ir_i$.  Linear independence follows successively from the
  first nonzero coordinate of $\sum_i z_ir_i$.
\end{proof}

The local flooring error is

\begin{equation}
  e_i=Ba_{i+1}-Aa_i,\qquad -B<e_i<A.
  \label{eq:local-error}
\end{equation}

Using~\eqref{eq:qv-definition},

\begin{equation}
  Av_i-Bv_{i+1}=-B^m e_i.
  \label{eq:carry-error}
\end{equation}

Thus an adjacent carry changes the actual gap sum by exactly $-e_i$.

\subsection{From a real interior point to integer coefficients}

Fix a main mass $Q$.  The nonnegative real fiber

\[
  \mathcal P_Q=\{c\in\R_{\ge0}^m:q\cdot c=Q\}
\]

has a one-dimensional error image:

\begin{equation}
  \left\{\frac{v\cdot c}{Q}:c\in\mathcal P_Q\right\}
  =[\sigma_-,\sigma_+],
  \quad
  \sigma_-:=\min_i\sigma_i,
  \quad
  \sigma_+:=\max_i\sigma_i.
  \label{eq:slope-convex-hull}
\end{equation}

Indeed, $(v\cdot c)/Q$ is the convex combination of the $\sigma_i$ with
weights $q_ic_i/Q$.

\begin{lemma}[Interior representation]
  \label{lem:interior}
  Suppose $\tau$ lies at distance at least $0<\eta\le1$ from both endpoints of
  the interval in~\eqref{eq:slope-convex-hull}.  For any margin $S>0$, if
  \begin{equation}
    Q\ge C\frac{S}{\eta}\sum_iq_i,
    \label{eq:interior-mass}
  \end{equation}
  where $C$ depends only on the uniform slope bound, then there is a real
  vector $c^*$ with
  \[
    q\cdot c^*=Q,\qquad v\cdot c^*=\tau Q,\qquad c_i^*\ge S.
  \]
\end{lemma}

\begin{proof}
  Reserve $S$ units in every coordinate and write $c^*=S\mathbf1+y$.  Let
  $W=S\sum_iq_i$ and $Z=S\sum_iv_i$.  The reserved mass has slope
  $Z/W\in[\sigma_-,\sigma_+]$.  The remaining mass $Q-W$ must have slope
  \[
    \tau'=\frac{\tau Q-Z}{Q-W}.
  \]
  Its displacement from $\tau$ is at most a constant times $W/(Q-W)$, which
  is at most $\eta/2$ under~\eqref{eq:interior-mass} after increasing $C$.
  Hence $\tau'$ remains in the slope interval.  Distribute the remaining mass
  between coordinates attaining $\sigma_-$ and $\sigma_+$.
\end{proof}

The lemma is unchanged if one coordinate is duplicated with the same main
weight and a new uniformly bounded slope.  In that augmented fiber, the sum in
\eqref{eq:interior-mass} runs over all coordinates; one duplicate makes it at
most $2\sum_iq_i$.  This is the version used for a unit companion below.

\begin{lemma}[Local carry rounding]
  \label{lem:rounding}
  Let $c^{(0)}$ be an integer point and $c^*$ a real point with the same main
  coordinate.  Write, uniquely over $\R$,
  \[
    c^*-c^{(0)}=\sum_{i=1}^{m-1}z_i^*r_i.
  \]
  Replacing every $z_i^*$ by its nearest integer gives an integer point
  $\widehat c$ with the same main coordinate and
  \begin{equation}
    \|\widehat c-c^*\|_\infty\le(A+B)/2.
    \label{eq:local-rounding}
  \end{equation}
  If $c^*$ represents $kd$ in actual gap value, then
  \begin{equation}
    \left|\sum_i\widehat c_i a_i-kd\right|
    \le\frac12\sum_{i=1}^{m-1}|e_i|
    =\bigO_{A,B}(m).
    \label{eq:rounding-residual}
  \end{equation}
\end{lemma}

\begin{proof}
  The integer basis in \cref{prop:kernel} also spans the real kernel, so the
  displayed real decomposition exists and is unique.
  Each coordinate is affected by at most two adjacent carry-rounding errors,
  each of magnitude at most one half, proving~\eqref{eq:local-rounding}.
  Equation~\eqref{eq:carry-error} says that rounding carry $i$ changes actual
  value by $-e_i$ times its rounding error.  Summing proves
  \eqref{eq:rounding-residual}.
\end{proof}

Without an additional arithmetic property, the residual in
\eqref{eq:rounding-residual} need not be removable by ordinary carries.  The
unit companions are designed precisely for this obstruction.

\subsection{The unit direction and a conductor bound}

Suppose one future position $j$ contains both $a_j$ and $a_j+1$.  Give the
companion a separate coefficient.  The two coordinates have the same main
weight $q_j$, while their error weights differ by $B^m$.  Their slopes differ
by

\begin{equation}
  \frac{B^m}{q_j}=R^{-j}.
  \label{eq:unit-slope-width}
\end{equation}

More importantly, the exchange vector

\begin{equation}
  s=-\mathbf e_{a_j}+\mathbf e_{a_j+1}
  \label{eq:unit-direction}
\end{equation}

preserves the main coordinate and changes the actual gap sum by exactly one.

\begin{proposition}[Augmented kernel]
  \label{prop:augmented-kernel}
  If $q^+$ duplicates the weight $q_j$ for the companion coordinate, then
  \[
    \ker_{\Z}(q^+)
    =\left(\bigoplus_{i=1}^{m-1}\Z r_i\right)\oplus\Z s.
  \]
\end{proposition}

\begin{proof}
  Merge the companion coefficient into the original $j$th coefficient.  Any
  augmented kernel vector then becomes a vector in $\ker_{\Z}(q)$, so
  \cref{prop:kernel} expresses it in the $r_i$.  Restoring the companion
  coordinate leaves an integer multiple of $s$.  Directness follows because
  every $r_i$ has zero companion coordinate whereas $s$ has companion
  coordinate one.
\end{proof}

\begin{lemma}[Augmented rounding and exact correction]
  \label{lem:augmented-rounding}
  Let $c^{(0)}\in\Z_{\ge0}^{m+1}$ and
  $c^*\in\R_{\ge0}^{m+1}$ have the same augmented main coordinate.  Suppose
  $c^*$ represents $kd$ in actual gap value and every coordinate of $c^*$ is
  at least $S$.  There is a constant $C_{A,B}$ such that, whenever
  $S\ge C_{A,B}m$, the same main coordinate has a nonnegative integer
  representation of exactly $kd$.
\end{lemma}

\begin{proof}
  By \cref{prop:augmented-kernel}, over $\R$ there are unique coefficients
  $z_1^*,\ldots,z_{m-1}^*,z_u^*$ with
  \[
    c^*-c^{(0)}=\sum_{i=1}^{m-1}z_i^*r_i+z_u^*s.
  \]
  Round every coefficient to the nearest integer.  Ordinary carry rounding
  changes any coordinate from $c^*$ by at most $(A+B)/2$, and rounding the
  unit coefficient adds at most $1/2$ on the two unit-pair coordinates.  The
  resulting vector $\widetilde c$ is integral and has the same main
  coordinate.  Its actual-value residual
  \[
    r:=\sum_i\widetilde c_i a_i+\widetilde c_u(a_j+1)-kd
  \]
  is an integer and satisfies
  \[
    |r|\le\frac12\sum_{i=1}^{m-1}|e_i|+\frac12
         =O_{A,B}(m).
  \]
  Moving $-r$ integral steps along $s$ keeps the main coordinate and removes
  the residual exactly.  The total decrease of either unit-pair coefficient,
  as well as every ordinary rounding decrease, is at most $C_{A,B}m$ after
  fixing the constant.  The margin hypothesis therefore keeps all final
  coefficients nonnegative.
\end{proof}

\begin{theorem}[Unit-companion conductor]
  \label{thm:conductor}
  There are constants $t_*$ and $C_{A,B}$ with the following property.  Let
  $t\ge t_*$, let $d\in\{h_t,h_t+1\}$ be a gap present in the completed
  sequence, set $m=m(t)$, and put $a_i=h_{t+i}$ for $1\le i\le m$.  If the
  augmented future list contains both $a_j$ and $a_j+1$ for some
  $1\le j\le m$, then its multiplier semigroup $M_d^+$ has conductor
  \begin{equation}
    \cond M_d^+
    \le C_{A,B}\,mR^j\left(\frac d{B^m}+1\right).
    \label{eq:conductor-bound}
  \end{equation}
\end{theorem}

\begin{proof}
  Fix a multiplier $k$.  We construct nonnegative integer coefficients
  representing $kd$ whenever $k$ exceeds the right-hand scale.

  Choose a main coordinate of the form
  \[
    Q=kB^m-D=AX,
  \]
  with integer $X$.  The required error slope is
  \begin{equation}
    \tau(X)=\frac{Dd}{Q}
           =\frac{dkB^m}{AX}-d,
    \qquad
    X(\tau)=\frac{dkB^m}{A(d+\tau)}.
    \label{eq:slope-inverse}
  \end{equation}
  Because $t\ge t_*$ and every augmented slope has absolute value below $2$,
  $d+\tau\in[d-2,d+2]=\Theta(d)$.  Therefore, on the relevant slope interval,
  \[
    \left|\frac{dX}{d\tau}\right|
    =\Theta_{A,B}(kB^m/d).
  \]
  By~\eqref{eq:unit-slope-width}, the two unit-pair slopes delimit an interval
  of width $R^{-j}$.  Remove the outer thirds.  Every point of the remaining
  interval is at distance at least $R^{-j}/3$ from both endpoints of the full
  augmented slope interval.  Its preimage under
  \eqref{eq:slope-inverse} has length
  \[
    \Omega_{A,B}\!\left((kB^m/d)R^{-j}\right).
  \]
  For $k\ge C R^j(d/B^m+1)$ this preimage has length greater than one and
  therefore contains an integer $X$.  Uniformly on it,
  $X=\Theta_{A,B}(kB^m)$.  Since $A^{m-1}<h_t\le A^m$ and
  $F_m=O_{A,B}(A^{m-1})$, increasing the fixed constant $C_{A,B}$ also makes
  $X>F_m$.  Thus \cref{lem:compound} gives a nonnegative integer reference
  representation of $X$ in the $w_i$, equivalently of $Q$ in the $q_i$;
  extend it by a zero companion coefficient.

  The selected slope remains at distance
  $\eta=\Theta(R^{-j})$ from the full augmented interval endpoints.  The
  augmented main-weight sum is at most $2\sum_iq_i=O_{A,B}(A^m)$, so the
  augmented version of \cref{lem:interior} supplies an exact real
  representation with margin $S$ once
  \[
    k\ge C_{A,B}S R^jR^m.
  \]
  Using $R^m=\Theta(d/B^m+1)$ and taking $S=C_{A,B}m$ yields the scale in
  \eqref{eq:conductor-bound}.

  The reference point and this real point have the same main coordinate, and
  the real point satisfies both equations in
  \eqref{eq:coupled-coordinates}.  With $S=C_{A,B}m$,
  \cref{lem:augmented-rounding} converts it to a nonnegative integer point of
  exactly the same actual value $kd$.

  Thus all $k$ above~\eqref{eq:conductor-bound} are representable.
\end{proof}

The augmented list used in the theorem is a subset of the gaps already
executed in the actual Shellsort pass.  Adding any other executed gaps can only
enlarge the representable multiplier semigroup.  Therefore the same conductor
bound is a valid upper bound for the actual missing-multiplier count $U_d$.

The unit pair removes every residue-class obstruction constructively; no gcd
assumption on the future gaps is needed.

\subsection{Zero-density placement covers every window}

We first verify the guard without floating-point arithmetic.  Writing
$\alpha=p/q$ in lowest terms, the two comparisons are checked from
\[
  pA^{2723}<(10^{1000}+1)qB^{2723},\qquad
  pA^{2724}\ge(10^{1000}+1)qB^{2724}.
\]
These exact integer inequalities are respectively equivalent to

\begin{equation}
  h_{2723}\le10^{1000}<h_{2724}.
  \label{eq:guard-check}
\end{equation}

In particular every added companion is greater than $10^{1000}$.  To prove
window coverage, choose $r$ with $s_r\le t<s_{r+1}$ and use $s_{r+1}$ as the
next companion index (also when $t=s_r$).  Its offset satisfies
\[
  j=s_{r+1}-t
  \le s_{r+1}-s_r
  =\lceil\log(s_r+2)\rceil
  \le\lceil\log(t+2)\rceil.
\]
Since $m(t)=\kappa t+O(1)$ and $\kappa>0$, this logarithmic offset is smaller
than $m(t)$ for every sufficiently large $t$.  Thus

\begin{equation}
  1\le j\le\lceil\log(t+2)\rceil<m(t).
  \label{eq:nearby-companion}
\end{equation}

Hence every sufficiently large complete future window contains a unit pair.
The strict inequality $j<m(t)$ is also needed at the execution boundary: a
complete backbone window contains $h_{t+j+1}<N$, and our choice of $t_*$ gives
$h_{t+j}+1<h_{t+j+1}$.  Thus both members of the unit pair are below $N$ and
have actually been executed before the current pass.
Moreover $t=\Theta(\log d)$, so

\begin{equation}
  R^j\le R(t+2)^{\log R}=\polylog d.
  \label{eq:unit-polylog}
\end{equation}

Combining \cref{thm:conductor}, \eqref{eq:window-scales}, and
$m=\bigO(\log d)$ yields

\begin{equation}
  U_d=\bigO(d^\kappa\polylog d)
  \label{eq:hole-bound}
\end{equation}

for every sufficiently large complete-window original or companion pass.

It remains to count the companions.  Let $T=\Theta(\log N)$ be the largest
backbone index below $N$.  The recurrence~\eqref{eq:companion-indices} advances
by $\Theta(\log s)$ at index $s$.  In an index block
$[2^\ell,2^{\ell+1})$, every step has length at least $c\ell$, so the block
contains $O(2^\ell/\ell)$ companion indices.  Summing these bounds through the
last block is dominated, up to a constant, by $O(T/\log T)$.  Therefore

\begin{equation}
  \#\{r:h_{s_r}<N\}
  =\bigO\!\left(\frac{\log N}{\log\log N}\right)
  =o(\log N).
  \label{eq:companion-count}
\end{equation}

\subsection{Summing all passes}

Substituting~\eqref{eq:hole-bound} into the combined pass bound gives, for an
internal pass,

\begin{equation}
  \cost_d(N)=\bigO\!\left(
    \min\left\{\frac{N^2}{d},\ Nd^\kappa\polylog d\right\}
  \right).
  \label{eq:internal-pass}
\end{equation}

Ignoring polylogarithmic factors, the two terms balance when

\[
  d^{1+\kappa}\asymp N.
\]

At that scale their common exponent is

\begin{equation}
  \beta
  =2-\frac1{1+\kappa}
  =1+\frac{\kappa}{1+\kappa}
  =1.02429645165747606869\ldots.
  \label{eq:final-exponent}
\end{equation}

Set $D=N^{1/(1+\kappa)}$.  Because $h_t=\Theta(R^t)$ and every companion is
only $h_t+1$, there are at most two gaps at each geometric scale.  Hence

\begin{equation}
  \sum_{d\le D}d^\kappa=O(D^\kappa),
  \qquad
  \sum_{d>D}\frac1d=O(D^{-1}).
  \label{eq:appendix-geometric-sums}
\end{equation}

For complete-window passes with $d\le D$, use the representation term in
\eqref{eq:internal-pass} and the first sum above.  Their total is
$O(ND^\kappa\polylog N)$.  For complete-window passes with $d>D$, use the
insertion term and the second sum; their total is $O(N^2/D)$.

The largest gaps have fewer than $m(t)$ future backbone gaps below $N$.  Let
$h_T<N$ be the largest backbone gap used.  An incomplete window satisfies
$t+m(t)>T$.  From \eqref{eq:window-index},

\[
  t>\frac{T}{1+\kappa}-O(1),
  \qquad
  h_t=\Omega\!\left(N^{1/(1+\kappa)}\right).
\]

The same lower bound holds for a companion $h_t+1$.  Thus every incomplete
window is on the $d=\Omega(D)$ side.  Using $N^2/d$ and the second sum in
\eqref{eq:appendix-geometric-sums} bounds all top passes together by

\[
  \bigO\!\left(
    N^{2-1/(1+\kappa)}
  \right)
  =\bigO(N^\beta).
\]

Finally, every fixed early gap has $O(N)$ cost once the first later unit pair is
available: the two consecutive fixed integers generate a numerical semigroup
with fixed conductor.  Before that pair enters the gap list, $N$ lies in a
fixed finite interval and is absorbed into the asymptotic constant.  Finitely
many early passes therefore contribute only $O(N)$.

Since
\[
  ND^\kappa=N^{1+\kappa/(1+\kappa)}=N^\beta,
  \qquad
  N^2/D=N^{2-1/(1+\kappa)}=N^\beta,
\]
all pass classes total $O(N^\beta\polylog N)$.  This proves part 3 of
\cref{thm:main} without assuming a full future window at the top boundary.

Part 1 follows from~\eqref{eq:guard-check}: Shellsort uses only gaps below $N$,
and every companion exceeds $10^{1000}$.  Part 2 follows from the geometric
backbone and~\eqref{eq:companion-count}.  For part 4, every backbone term lies
within one of $\alpha(A/B)^t$, every companion lies within one of the same
quantity, and a single larger fixed constant covers the finite tuned prefix.
Thus Zang's theorem applies with $a=A$, $b=B$, and $q=\alpha$.  If
$\gamma=\log_B A$, its exponent is
\[
  1+\frac{\gamma-1}{2\gamma-1}
  =1+\frac{\kappa}{1+\kappa}
  =1.02429645165747606869\ldots,
\]
where $\kappa=1-1/\gamma$.  A swap lower bound is also an operation-count
lower bound.  This completes all four parts of the theorem.

\subsection{Why finite-prefix tuning preserves the theorem}

The proof above used the geometric formula only from a sufficiently large
index onward.  Replacing the eight positive backbone terms below 1416 by
$P_8$ changes finitely many passes.  After unit companions appear, each fixed
replacement gap again has a fixed multiplier-semigroup conductor and hence
$\bigO(N)$ pass cost.  A finite sum of such terms cannot change the exponent in
\eqref{eq:main-bound}.  This establishes the stated theorem specifically for
$\mathcal H^\sharp_{\rm tune}$, rather than only for the pre-tuning backbone.

For completeness, a sparser alternative places the next companion after
$\Theta(m(s_r))$ indices.  It adds only $\bigO(\log\log N)$ gaps, but the factor
$R^j$ in~\eqref{eq:conductor-bound} then contributes a second $d^\kappa$ and
gives the weaker exponent

\[
  1+\frac{2\kappa}{1+2\kappa}
  =1.04744027301502511325\ldots.
\]

The completion in \cref{thm:main} chooses the denser, still zero-density
placement because it gives the strongest exponent while remaining invisible
through the guard threshold.
